\chapter{Validación y resultados}

\section{Estrategia de validación}

La validación de Nuvlo combina pruebas automatizadas, revisión manual de la interfaz y comprobación con datos reales de Jira. Esta estrategia busca cubrir dos aspectos distintos: que los cálculos internos sean correctos y que la aplicación sea usable en los flujos principales.

Para conectar implementación y requisitos se utiliza una matriz requisito-prueba. En ella, cada requisito funcional, no funcional o de integración se asocia con una prueba, una revisión manual o una evidencia de uso. La matriz completa se recoge en la Tabla~\ref{tab:matriz-requisito-prueba}, dentro del capítulo de requisitos, y se resume en este capítulo desde el punto de vista de la validación.

\section{Tipos de prueba utilizados}

Antes de interpretar la matriz requisito-prueba conviene distinguir los tipos de comprobación empleados durante el desarrollo:

\begin{itemize}
    \item \textbf{Pruebas unitarias}: verifican funciones aisladas, especialmente cálculos de métricas, percentiles, validaciones de datos y reglas de alerta. Se ejecutan con Vitest~\cite{VitestGuide}.
    \item \textbf{Pruebas de integración}: comprueban que varias piezas funcionan juntas, por ejemplo API, repositorios, PostgreSQL, Redis y servicios de sincronización.
    \item \textbf{Pruebas con Jira real}: validan el flujo OAuth, la lectura de proyectos, la sincronización de incidencias y el almacenamiento local de datos procedentes de Jira Cloud.
    \item \textbf{Pruebas con dataset offline}: usan datos controlados para comprobar cálculos que requieren histórico suficiente, como Lead Time, Cycle Time, percentiles o Velocity.
    \item \textbf{Pruebas end-to-end}: recorren la aplicación desde el navegador, incluyendo navegación, filtros, widgets, alertas, logs y modo demo. Se plantean con Playwright~\cite{PlaywrightDocs}.
    \item \textbf{Revisión visual y responsive}: comprueba capturas en escritorio y móvil para validar legibilidad, navegación y coherencia con las vistas previstas.
\end{itemize}

\section{Pruebas automatizadas}

Las pruebas unitarias se centran en cálculos de métricas, percentiles, validaciones y reglas de alerta. Para ello se utiliza Vitest~\cite{VitestGuide}. Las pruebas de integración comprueban el comportamiento de la API, la persistencia y los servicios que dependen de PostgreSQL o Redis. En los casos relacionados con Jira se utilizan datos controlados o respuestas simuladas para evitar depender siempre de una conexión externa.

Las pruebas end-to-end se plantean con Playwright~\cite{PlaywrightDocs}. Su objetivo es recorrer los flujos principales desde el punto de vista del usuario: entrada en modo demo, navegación por vistas, uso de filtros, sincronización, consulta de alertas y revisión del registro de actividad.

\section{Matriz requisito-prueba}

La matriz completa de validación se muestra en la Tabla~\ref{tab:matriz-requisito-prueba}. Para evitar repetir en este capítulo toda la tabla, la Tabla~\ref{tab:resumen-validacion-requisitos} resume la lectura global de esa matriz.

{\renewcommand{\arraystretch}{1.22}\footnotesize
\begin{longtable}{p{3.0cm}p{3.1cm}p{6.0cm}}
\caption{Resumen de la matriz requisito-prueba}\label{tab:resumen-validacion-requisitos}\\
\toprule
\textbf{Bloque} & \textbf{Estado general} & \textbf{Evidencia principal} \\
\midrule
\endfirsthead
\toprule
\textbf{Bloque} & \textbf{Estado general} & \textbf{Evidencia principal} \\
\midrule
\endhead
Integración Jira y autenticación & Cumplido & Flujo OAuth 2.0 3LO, sesión protegida, sincronización de proyecto real y persistencia en PostgreSQL. \\
Dashboard y visualización & Cumplido & Dashboard responsive, filtros rápidos, widgets configurables, ayudas contextuales y capturas de escritorio/móvil. \\
Métricas principales & Cumplido & Cálculo de WIP, Velocity, Lead Time, Cycle Time, media y percentiles con dataset offline y datos sincronizados. \\
Alertas y actividad & Cumplido & Reglas de umbral, contador visual de alertas, notificación emergente, historial de alertas y registro de actividad. \\
Retención, seguridad y mantenimiento & Cumplido & Cookies \texttt{httpOnly}, CSRF, PKCE, cifrado de tokens, Redis con TTL y limpieza de logs/alertas antiguos. \\
Configuración avanzada y métricas futuras & Parcial/Futuro & Configuración básica disponible; quedan como ampliación la configuración avanzada completa, Deployment Frequency y más métricas de flujo. \\
\bottomrule
\end{longtable}
\vspace{0.45em}
}


\section{Validación con Jira real}

La aplicación se ha probado con un proyecto real de Jira creado para demostración. Este escenario permite validar el flujo OAuth, la obtención de proyectos accesibles, la sincronización de incidencias y la persistencia en la base de datos local. Las capturas del manual de usuario muestran la aplicación funcionando con datos procedentes de Jira.

La validación con Jira real también sirve para detectar limitaciones prácticas. Por ejemplo, si las incidencias se crean y cierran en un intervalo muy corto, Lead Time y Cycle Time pueden aparecer como cero o valores poco representativos. Esto no implica que el cálculo esté ausente, sino que el histórico disponible no contiene todavía información temporal suficiente para un análisis realista.

\section{Validación con demo offline}

La demo offline cubre el caso contrario: proporciona datos preparados para que todas las métricas tengan valores visibles y coherentes. Este modo es útil para presentaciones, pruebas sin conexión y comprobaciones de regresión visual. Al no depender de la API de Jira, permite demostrar el comportamiento esperado de la interfaz incluso si la autenticación externa falla o no hay red disponible.

\section{Resultados obtenidos}

El resultado actual es un MVP funcional de Nuvlo. La aplicación permite iniciar sesión con Atlassian, sincronizar información de Jira, consultar métricas ágiles, aplicar filtros, configurar widgets, crear alertas y revisar actividad. Además, dispone de una demo offline que actúa como respaldo para validación y presentación.

El sistema todavía no pretende cubrir todos los escenarios de una herramienta comercial. Algunas mejoras, como métricas avanzadas, despliegue productivo, integración con repositorios de código o análisis predictivo, quedan fuera del alcance principal y se recogen como líneas de trabajo futuro.
